\chapter{Méthodologie et résultats}

La méthodologie décrite dans ce chapitre expose l’ensemble des étapes permettant de constituer, d’aligner et
d’exploiter un corpus contenant des textes en français, alsacien, corse et occitan. Le but est d’entraîner
un modèle capable d’identifier automatiquement des
paires de phrases parallèles. Pour ce faire nous utiliserons des outils d’alignement interlinguistique,
des techniques de représentation sémantique, ainsi que des méthodes statistiques de classification supervisée.

Dans un premier temps, nous utiliserons l’outil \textit{InterText} afin d'aligner les données issues de
l’Atlas des langues régionales du LIMSI. Ce premier travail nous permet d'avoir une autre partie de données pour le corpus,
avec PARCOLAB et la parabole de l'enfant prodigue. Nous calculerons ensuite des scores de similarité textuelle
sur ce corpus.

Les similarités sémantiques issues du modèle SBERT seront
exploitées pour attribuer des scores aux comparaisons interphrastiques sur des données parallèles,
semi-parallèles ou non parallèles issues du corpus d'entraînement. Ces scores, combinés à d'autres variables
complémentaires mesurant la proximité structurelle et phonétique entre les phrases,
serviront de variables d'entrée pour l'entraînement d'un modèle de régression logistique capable
d'identifier automatiquement les phrases parallèles.

Enfin, nous présentons une étape complémentaire permettant d’évaluer le bon fonctionnement du modèle en nous
appuyant sur le corpus établi à partir du programme de \citet{tsolakis_2024}. Après validation du modèle sur ce corpus de
référence, il nous faudra étendre le code du modèle afin d'identifier non plus seulement des phrases
parallèles, mais également des fragments semi-parallèles.

\newpage

\section{InterText}

\textit{InterText}\footnote{\hyperlink{https://wanthalf.saga.cz/intertext}{InterText}} \citep{vondricka_2014} est un environnement conçu pour faciliter
l’annotation et l’alignement de textes multilingues. Développé à l’origine pour les études contrastives et les
sciences du texte, \textit{InterText} permet de comparer des documents appartenant à des langues différentes
et de produire des alignements fiables au niveau phrase ou paragraphe. Pour réaliser ces alignements \textit{InterText}
se base sur des aligneurs déjà existant
telles que Hunalign \citep{nicolov_2007}, qui exploite des heuristiques basées sur la longueur des phrases et la
fréquence des mots et TCA2 \citep{hofland_1996, hofland_1998, Vondricka_2010}, un aligneur développé
initialement pour des corpus de langues scandinaves.

\begin{figure}[h]
    \centering
    \includegraphics[width=\textwidth]{ui_intertext.png}
    \caption{L'interface utilisateur d'InterText.}
    \label{ui_intertext}
\end{figure}

Son principal avantage réside dans sa capacité à fournir une interface (voir \ref{ui_intertext}) permettant
d’afficher côte à côte deux versions d’un même document (ou de deux documents comparables) afin d’aligner
semi-automatiquement ou manuellement les segments. Concrètement, l'utilisateur
importe deux documents au format XML ou txt, qui sont ensuite segmentés automatiquement en unités textuelles
(phrases, syntagmes, ou paragraphes) selon des règles de segmentation paramétrables. Ces segments peuvent
être alignés de trois manières{ :}

\begin{itemize}
\item \textbf{Alignement automatique}{ :} l'outil s'appuie sur Hunalign ou TCA2 pour établir des
correspondances probabilistes entre segments sur la base d'indices lexicaux ou structurels.
\item \textbf{Alignement assisté}{ :} l'utilisateur valide ou corrige les propositions générées
automatiquement. Cette modalité constitue un compromis intéressant entre efficacité et précision,
permettant de corriger les erreurs de l'alignement automatique tout en réduisant la charge de travail manuel.
\item \textbf{Alignement manuel}{ :} utilisé lorsque la langue ou la variation dialectale rendent
les procédures automatiques peu fiables. C'est cette modalité qui a été principalement mobilisée
dans notre cas, compte tenu de la forte variabilité dialectale des textes de l'Atlas LIMSI.
\end{itemize}

Au terme du processus, \textit{InterText} permet d'exporter ces
alignements sous forme de paires structurées ou de fichiers textes séparés. Pour un corpus 
tel que celui de l’Atlas LIMSI, qui n’a pas été initialement conçu comme un corpus parallèle,
\textit{InterText} représente un outil indispensable permettant de pallier l’absence d’alignement natif.

\section{Alignement du corpus LIMSI avec InterText}

Le corpus issu de l’Atlas des langues régionales du LIMSI est un corpus comparable. En effet, bien que tous
les textes reposent sur une traduction de la fable
\textit{La bise et le soleil} d'Ésope, les variations dialectales, les découpages phrastiques
et les divergences d’interprétation rendent l’alignement automatique impossible.

Dans cette perspective, \textit{InterText} a été mobilisé afin de constituer un corpus parallélisé entre les
variétés régionales. L’alignement s’effectue à partir du texte de référence en français et des versions du
poème dans les langues régionales de France sur lesquelles nous travaillons.
Pour chaque paire de documents, les segments sont alignés un à un lorsque la correspondance est claire,
ou fusionnés lorsqu'une phrase française correspond à plusieurs segments dans la langue régionale, ou
inversement.

Cette étape produit un ensemble de paires alignées, accompagné d’un nombre de
segments imparfaits non négligeables{ :} phrases plus longues d’un côté que de l’autre,
segments fusionnés ou fragmentés, etc.
Ce phénomène justifie la nécessité d'introduire plus loin une notion de \textit{semi-parallélisme}.
Cette notion permettra au modèle de traiter ces cas intermédiaires sans les rejeter systématiquement
comme des paires non parallèles.

\section{Séparation du corpus}

%CORPUS_DIR = "corpus/"
Dans un premier temps, le programme procède à la constitution d'un ensemble de paires de fichiers textes, 
chacune associant un document source en français à son équivalent dans une langue régionale. 
Cette organisation repose sur une logique propre à chaque corpus. Pour l'Atlas LIMSI, le texte français de 
référence est systématiquement mis en regard de chacune des versions dialectales disponibles. Pour la parabole 
du fils prodigue, le même principe est appliqué : le texte pivot en français est apparié avec chaque 
traduction régionale. Le corpus ParCoLab, quant à lui, requiert un traitement légèrement différent : 
les fichiers y étant nommés selon une convention incluant l'indication de langue, le programme identifie 
dynamiquement les fichiers français et leur associe les fichiers de la même collection dans les autres 
langues.

Chaque paire est ainsi définie par trois éléments : le fichier source en français, le fichier cible en 
langue régionale, et un identifiant permettant de tracer l'origine de la paire au sein du corpus. Cette 
identification est importante, car elle permet, lors des analyses ultérieures, de distinguer les paires
selon leur provenance et donc selon leur profil linguistique et leur niveau de bruit attendu.

Dans un second temps, le programme procède au chargement effectif des fichiers. Chaque fichier est 
lu ligne par ligne, chaque ligne correspondant à une phrase alignée. Le programme vérifie que les deux 
fichiers d'une même paire contiennent le même nombre de phrases, condition nécessaire pour que l'alignement 
soit considéré comme valide. Si cette condition n'est pas remplie, la paire est écartée et un avertissement 
est émis. Cette vérification garantit la cohérence structurelle des données d'entraînement avant toute 
opération de calcul de similarité.

\subsection{Création de données semi et non-parallèles}

%def build_parallel_subset(pairs, size=None):
Une fois les paires parallèles chargées, le programme procède à la constitution de trois sous-ensembles 
distincts destinés à l'entraînement du modèle de régression logistique : un sous-ensemble de paires 
parallèles, un sous-ensemble de paires semi-parallèles, et un sous-ensemble de paires non parallèles.

Le sous-ensemble parallèle est constitué directement à partir des paires chargées précédemment. Si une 
taille maximale a été définie par l'utilisateur, un échantillonnage aléatoire est effectué afin de ne 
retenir qu'un nombre limité de paires. Dans le cas contraire, l'intégralité des paires disponibles est 
conservée.

Le sous-ensemble semi-parallèle est généré artificiellement à partir de ces mêmes paires parallèles, en 
tronquant le segment en langue régionale. Deux stratégies de troncature sont appliquées selon la structure 
de la phrase. La stratégie privilégiée consiste à identifier la première ponctuation interne de la phrase 
(virgule, point-virgule ou deux-points) et à ne conserver que la portion qui la précède. Cette approche est 
motivée par le fait que les transcriptions dialectales de l'Atlas LIMSI présentent fréquemment des omissions 
ou des condensations en milieu ou en fin de phrase, rendant cette troncature linguistiquement plausible. 
Lorsqu'aucune ponctuation interne n'est disponible, une stratégie de repli consiste à ne conserver que la 
première moitié des mots de la phrase. Dans les deux cas, le segment français d'origine est conservé intact, 
de sorte que la paire résultante présente une correspondance partielle entre les deux langues.

Le sous-ensemble non parallèle est quant à lui construit par appariement aléatoire de phrases françaises et 
de phrases en langue régionale provenant de positions différentes dans le corpus. Le programme vérifie 
systématiquement que les paires ainsi générées ne figurent pas déjà parmi les paires parallèles afin d'éviter 
l'introduction de faux négatifs dans les données d'entraînement. Dans le cas de corpus de petite taille, le 
nombre de combinaisons distinctes disponibles peut être insuffisant pour atteindre la taille cible, ce qui est 
signalé à l'utilisateur.

Les trois sous-ensembles sont enfin fusionnés et mélangés aléatoirement afin de constituer le jeu de 
données complet utilisé pour l'entraînement du classifieur.

\section{Calcul des similarités sémantiques avec SBERT}

Une fois l’alignement réalisé, chaque paire de phrases est comparée au moyen d’un modèle reposant sur
l’architecture SBERT\citep{reimers_2019}. SBERT génère pour chaque phrase un embedding, obtenu
par agrégation des sorties de l'encodeur via un mécanisme de \textit{mean pooling}. La similarité entre
deux phrases est ensuite calculée par similarité cosinus entre leurs embeddings respectifs,
produisant un score compris entre 0 et 1.

Comme modèle SBERT nous décidons d'utiliser le modèle 
\textit{paraphrase-multilingual-MiniLM-L12-v2}\footnote{\hyperlink{https://huggingface.co/sentence-transformers/paraphrase-multilingual-MiniLM-L12-v2}{paraphrase-multilingual-MiniLM-L12-v2}} \citep{reimers_2019}.
MiniLM est un modèle distillé, il a été entraîné à imiter le comportement d'un modèle plus grand en 
reproduisant ses relations d'attention. Le suffixe L12 indique qu'il possède 12 couches de
transformer, et la dimension de ses embeddings est de 384 dimensions. Un autre modèle nommé 
\textit{paraphrase-multilingual-mpnet-base-v2}\footnote{\hyperlink{https://huggingface.co/sentence-transformers/paraphrase-multilingual-mpnet-base-v2}{paraphrase-multilingual-mpnet-base-v2}} \citep{reimers_2019}
pourrait aussi être intéressant. Le tableau ci-dessous présente les caractéristiques comparatives des deux modèles{ :}
\begin{table}[h!]
\centering
\begin{tabular}{lccc}
\hline
\textbf{Caractéristique} & \textbf{mpnet-base-v2} & \textbf{MiniLM-L12-v2} \\
\hline
Dimension des embeddings & 768 & 384 \\
Nombre de paramètres & ~278M & ~118M \\
Vitesse d'inférence & Modérée & Rapide \\
Qualité sémantique & Supérieure & Inférieure \\
Mémoire nécessaire & Élevée & Faible \\
\hline
\end{tabular}
\caption{Tableau comparatif des modèles MPNet et MiniLM}
\end{table}
Le choix de MiniLM se justifie pour plusieurs raisons complémentaires. Tout d'abord,
les embeddings à 384 dimensions restent suffisamment riches pour capturer les similarités sémantiques
entre phrases françaises et régionales. De plus la légère perte de qualité de MiniLM par rapport à MPNet
aura peu d'impact visible sur les résultats finaux étant donné la taille modeste de nos corpus.
Enfin, la réduction de la mémoire nécessaire et l'accélération de l'inférence deviennent des avantages
concrets lors de l'étape de fragmentation, où le nombre de paires à comparer croît significativement.
Les deux modèles étant entraînés sur des données de paraphrase dans plus de 50 langues,
ils sont tous deux capables de projeter des phrases de langues différentes dans un espace sémantique partagé.
C'est indispensable pour comparer directement le français aux langues régionales sans traduction
intermédiaire.

Pour entraîner un modèle supervisé à reconnaître les phrases parallèles, il est nécessaire de disposer de
données annotées selon différents degrés de correspondance. Nous distinguons trois catégories{ :}

\begin{itemize}
    \item Parallèles{ :}
    \[
    \text{phrase A} \Rightarrow \text{phrase A traduite.}
    \]
    Les deux segments expriment exactement la même information, même si la structure syntaxique varie.
    
    \item Semi-parallèles{ :}
    \[
    \text{phrase A} \Rightarrow \text{phrase A traduite, mais avec une partie manquante.}
    \]
    Ce cas se rencontre fréquemment lorsque certains détails sont omis ou condensés dans une version dialectale.
    Dans notre mémoire nous allons devoir créer artificiellement des variations de phrases existantes afin
    de produire des paires ayant un contenu informationnel similaire sans être des traductions exactes.
    Par exemple, une phrase peut être transformée par paraphrasage automatique, par suppression de certains
    éléments d’information (comme une suite de phrase après une virgule) ou par reformulation partielle.
    
    \item Non parallèles{ :}
    \[
    \text{phrase A} \Rightarrow \text{phrase B distincte sans corrélation sémantique.}
    \]
    Ces paires servent de négatifs dans l’apprentissage supervisé.
\end{itemize}

Pour chacune de ces catégories, un score STS est calculé à partir des
embeddings SBERT par similarité cosinus. Ces scores, accompagnés des variables complémentaires décrites ci-dessous,
constituent la base d’apprentissage du modèle de régression logistique.

\section{Principe du modèle de régression logistique}

La régression logistique est un modèle statistique utilisé pour prédire une probabilité appartenant à deux
catégories (Dans notre cas{ :} parallèle / non parallèle). La régression logistique transforme la sortie en une
probabilité comprise entre 0 et 1 grâce à la fonction sigmoïde{ :}
\[
\sigma(x) = \frac{1}{1 + e^{-x}}
\]

Elle constitue un modèle simple pour classifier les paires de phrases selon
qu’elles soient parallèles, semi-parallèles ou non parallèles. Afin d'enrichir la capacité discriminante du
modèle, plusieurs variables d'apprentissage sont combinées en complément des scores STS produits par SBERT.

Les scores issus du modèle SBERT constituent la variable principale. Ils capturent la proximité sémantique
globale entre deux phrases dans un espace vectoriel multilingue partagé. Cette méthode est particulièrement adaptée à
la comparaison entre le français et les langues régionales.

\subsection{Score de longueur}

Une variable complémentaire mesure la similarité structurelle entre deux phrases par le biais de leur
longueur respective. Ce score est calculé en divisant la valeur absolue de la différence du nombre de mots
des deux phrases par la longueur de la phrase la plus longue{ :}
\[ score de longueur = \frac{\left| len(P_{1})-len(P_{2}) \right|}{max(len(P_{1}),len(P_{2}))} \]
Où $P$ correspond à une phrase.
Un score proche de 0 indique des phrases de longueur similaire, ce qui constitue un indice favorable au
parallélisme. Cette mesure s'inscrit dans la tradition des premières méthodes d'alignement statistique,
comme celle de \citet{gale_1991}, qui exploitaient déjà la corrélation entre longueurs de phrases pour
établir des correspondances probabilistes.

\subsection{Indices de Jaccard}

Afin d'exploiter les similarités de surface entre les phrases, notamment les emprunts lexicaux et les
transfuges fréquents entre le français et les langues régionales, des indices de Jaccard
sont calculés sur les bigrammes, trigrammes et quadrigrammes de caractères. L'indice de Jaccard mesure le
chevauchement entre deux ensembles en divisant la taille de leur intersection par la taille de leur union.
Avant ce calcul, les phrases sont normalisées par suppression de la casse, de la ponctuation et des accents,
afin de réduire le bruit lié aux variations orthographiques, très fréquentes dans les langues
régionales dont la normalisation est parfois incomplète.

\subsection{Représentations pseudo-phonétiques}

Une dernière variable exploite les similarités phonétiques entre les phrases via la bibliothèque Python
jellyfish \citep{turk_2013}. Les phrases sont d'abord encodées selon l'algorithme Metaphone \citep{philips_1990},
qui représente les mots selon leur prononciation approximative de sorte que des mots phonétiquement proches
partagent une même représentation. Par exemple, les mots français {« }chant{ »} et {« }champ{ »} produiraient des
représentations Metaphone similaires. Un indice de Jaccard est ensuite calculé sur ces représentations
phonétiques. Cette variable est particulièrement pertinente dans notre contexte, car les langues régionales de
France partagent souvent avec le français des racines étymologiques communes dont la proximité est davantage
perceptible à l'oral qu'à l'écrit, notamment en occitan ou en corse. Elle permet ainsi de capturer des
correspondances que les mesures purement lexicales ou sémantiques pourraient manquer.

Le modèle apprend ainsi à combiner ces différentes variables pour reconnaître les configurations
typiques des phrases parallèles et à distinguer les paires non parallèles. L'entraînement sur des données
semi-parallèles devrait de plus rendre le modèle capable d'identifier ce type de segments, dont la présence
est particulièrement fréquente dans les corpus dialectaux comme celui de l'Atlas LIMSI.

\subsection{Application programmatique}

%def normalize(text):
Avant tout calcul, une étape de normalisation est appliquée aux textes. Celle-ci consiste à convertir les 
caractères en minuscules, à supprimer les accents par décomposition unicode, puis à retirer la ponctuation. 
Cette normalisation est particulièrement importante dans le contexte des langues régionales que nous étudions.
Celles-ci possèdent, en effet, des variations orthographiques dépendant du dialecte qui doivent être normalisées
sous peine de poser problème plus tard.

%def length_score(s1, s2):
Enfin, une vérification de cohérence est effectuée sur la première paire chargée (voir \ref{sanity_check}), 
affichant les valeurs de chacun de ces indicateurs afin de s'assurer du bon fonctionnement des fonctions 
avant le traitement de l'ensemble du corpus.

\begin{figure}[h]
    \centering
    \includegraphics[width=\textwidth]{sanity_check.png}
    \caption{Vérification de cohérence sur la première paires de phrases}
    \label{sanity_check}
\end{figure}

%print(f"Loading SBERT model: {SBERT_MODEL} ...")
Une fois les variables de surface calculées, le programme procède au calcul des scores de similarité
sémantique à l'aide du modèle SBERT. Le modèle est d'abord chargé en mémoire, puis appliqué à l'ensemble
des paires constituant le jeu de données.

Pour chaque paire de phrases, le modèle génère un plongement vectoriel pour chacun des deux segments,
c'est-à-dire une représentation numérique de la phrase dans un espace sémantique multilingue partagé.
La similarité entre les deux phrases est ensuite calculée par similarité cosinus entre leurs plongements
respectifs, produisant un score compris entre 0 et 1. Afin d'optimiser l'utilisation de la mémoire,
le calcul est effectué par lots plutôt que phrase par phrase. Les scores légèrement négatifs, qui peuvent
ponctuellement résulter du calcul de la similarité cosinus, sont ramenés à zéro afin de maintenir une
échelle de valeurs cohérente.

Ces scores constituent la variable principale du modèle de régression logistique, dans la mesure où ils
capturent la proximité sémantique globale entre les phrases françaises et leurs équivalents en langue
régionale dans un espace vectoriel commun, sans nécessiter de traduction intermédiaire entre les deux
langues.

%Computing complementary features
L'ensemble des scores calculés précédemment est ensuite rassemblé dans un tableau structuré, où chaque
ligne correspond à une paire de phrases et chaque colonne à une variable descriptive. Pour chaque paire,
le tableau enregistre les phrases source et cible, l'identifiant de la paire de langues, la catégorie de
parallélisme sous forme textuelle, ainsi qu'une étiquette binaire opposant les paires parallèles et
semi-parallèles d'un côté aux paires non parallèles de l'autre. Cette étiquette binaire constitue la
variable cible du modèle de régression logistique. À ces informations s'ajoutent l'ensemble des variables
d'apprentissage calculées précédemment : le score de similarité sémantique SBERT, le score de longueur,
les indices de Jaccard sur les bigrammes, trigrammes et quadrigrammes de caractères, ainsi que l'indice de
Jaccard phonétique. Ce tableau constitue la base d'apprentissage complète qui sera transmise au classifieur
lors de l'étape d'entraînement.

\newpage

\section{Données annotées}

\begin{figure}[h]
    \centering
    \includegraphics[width=\textwidth]{feature_distributions_kde.png}
    \caption{feature distributions kde}
    \label{feature_distributions_kde}
\end{figure}

Ce graphique présente les distributions de densité des six variables d'apprentissage selon les trois
catégories de parallélisme, sous forme de courbes KDE. Les tirets verticaux indiquent la médiane de chaque
classe. Il constitue un complément analytique essentiel à la matrice de corrélation et au graphique des
coefficients, en montrant non plus des valeurs agrégées mais la forme complète des distributions.

Score SBERT constitue la variable dont le pouvoir discriminant est le plus lisible. Les trois distributions
présentent des profils clairement distincts. La classe non parallèle est fortement concentrée dans les
valeurs basses avec un pic prononcé autour de 0,10 à 0,15, et une médiane proche de 0,18. La classe
semi-parallèle présente une distribution bimodale avec un premier mode autour de 0,15 et un second autour
de 0,40, reflétant l'hétérogénéité des paires tronquées : certaines troncatures importantes produisent des
scores faibles proches de ceux des paires non parallèles, tandis que les troncatures légères conservent une
proximité sémantique élevée. La classe parallèle se distingue par une distribution étalée vers les valeurs
hautes avec un mode principal autour de 0,60 à 0,90 et une médiane nettement supérieure aux deux autres
classes, aux alentours de 0,60. La superposition partielle entre les trois distributions, particulièrement
entre les classes semi-parallèle et non parallèle dans la région 0,10 à 0,30, explique une partie des
erreurs de classification observées dans les matrices de confusion.

Score de longueur présente un comportement atypique par rapport aux autres variables. Les trois classes
sont concentrées dans les valeurs basses, autour de 0 à 0,10, ce qui indique que la grande majorité des
phrases présentent des longueurs similaires quelle que soit leur catégorie. On observe néanmoins un pic
secondaire pour la classe semi-parallèle autour de 0,50, correspondant aux paires où la troncature à
mi-phrase a produit une asymétrie de longueur significative. La classe non parallèle présente une
distribution plus étalée avec une densité résiduelle jusqu'à 1,0, reflétant le fait que les phrases
appariées aléatoirement peuvent présenter des longueurs très différentes. Les médianes des trois classes
sont très proches les unes des autres, ce qui confirme le faible coefficient de cette variable dans le
modèle de régression logistique.

Jaccard Bigrammes montre une meilleure séparation que les variables de trigrammes et quadrigrammes. La
classe non parallèle est concentrée en 0 à 0,20 avec un pic autour de 0,15, tandis que la classe parallèle
présente une distribution nettement plus étalée avec un second mode visible autour de 0,40 à 0,50,
traduisant les chevauchements lexicaux fréquents entre le français et les langues régionales romanes.
La classe semi-parallèle occupe une position intermédiaire avec une distribution qui suit globalement
celle de la classe parallèle mais avec une densité plus faible dans les valeurs élevées.

Jaccard Trigrammes, Jaccard Quadrigrammes et Jaccard Phonétique présentent des distributions très
similaires entre elles, ce qui est cohérent avec la forte corrélation observée entre ces trois variables
dans la matrice de corrélation. Les trois classes sont massivement concentrées dans les valeurs très
basses, entre 0 et 0,05 à 0,10, avec des pics de densité extrêmement élevés en 0, ce qui explique les
valeurs d'axe Y atteignant respectivement 25, 80 et 20 unités de densité. Cette concentration reflète
le fait que les langues régionales, en particulier l'alsacien, partagent relativement peu de séquences
de caractères longues avec le français malgré leurs racines communes. La classe parallèle se distingue
par une queue de distribution légèrement plus étendue vers les valeurs positives et une médiane légèrement
supérieure, visible sur les tirets verts, mais la séparation entre classes reste faible pour ces trois
variables prises isolément. C'est leur combinaison avec le score SBERT au sein du modèle de régression
logistique, et non leur contribution individuelle, qui leur confère leur utilité discriminante, notamment
pour les paires où la proximité sémantique vectorielle est insuffisante à elle seule pour trancher.

\begin{figure}[h]
    \centering
    \includegraphics[width=\textwidth]{pie_by_language.png}
    \caption{Distribution des classes par paires de langues}
    \label{pie_by_language}
\end{figure}

Ces trois diagrammes circulaires présentent la distribution des classes au sein du jeu de données
d'entraînement, agrégée par paire de langues. Les effectifs indiqués entre parenthèses (fr-co : n=1186,
fr-gsw : n=834, fr-oc : n=1415) confirment le déséquilibre déjà observé dans l'analyse quantitative du
corpus : l'occitan est la langue la mieux représentée, suivi du corse puis de l'alsacien. Ce déséquilibre
est directement imputable à la surreprésentation des variétés occitanes dans l'Atlas LIMSI, dont les 47
textes dialectaux constituent la majorité des paires chargées.

Pour les trois paires de langues, la
distribution entre les trois classes est proche de l'équilibre tiers-tiers-tiers, ce qui est cohérent
avec la stratégie de construction du jeu de données décrite dans la méthodologie, où les sous-ensembles
parallèle, semi-parallèle et non parallèle sont construits à partir du même ensemble de paires parallèles
de base. On observe néanmoins de légères variations. Pour la paire fr-gsw, la classe non parallèle est
légèrement surreprésentée à 35,5\% au détriment de la classe parallèle à 32,3\%, ce qui peut s'expliquer
par les contraintes de génération des paires non parallèles dans de petits corpus : avec un nombre limité
de phrases disponibles, le programme génère plus facilement des paires négatives distinctes que dans
des corpus plus volumineux. Pour la paire fr-oc, la distribution est quasi parfaitement équilibrée à
34,3\% pour les classes parallèle et semi-parallèle, et 31,4\% pour la classe non parallèle, ce qui
s'explique par la taille plus importante de ce corpus permettant un échantillonnage plus équilibré.
Ces distributions globalement équilibrées constituent une condition favorable à l'entraînement du
classifieur, évitant les biais induits par des déséquilibres importants entre classes.

\begin{figure}[!h]
    \centering
    \includegraphics[width=\textwidth]{sbert_by_language_pair.png}
    \caption{Score SBERT moyen par paires de langues}
    \label{sbert_by_language_pair}
\end{figure}

Ce graphique en barres horizontales présente les scores SBERT moyens par paire de langues et par catégorie
de parallélisme, permettant une comparaison directe du pouvoir discriminant du modèle selon la langue
régionale considérée.

La hiérarchie attendue entre les trois classes est respectée pour l'ensemble des paires de langues : les
paires parallèles obtiennent systématiquement les scores les plus élevés, suivies des paires semi-parallèles,
puis des paires non parallèles. Ce résultat valide le fait que le modèle SBERT capture efficacement les
différences de parallélisme entre catégories indépendamment de la langue régionale cible.

La paire fr-co présente les scores les plus élevés pour les classes parallèle et semi-parallèle, atteignant
respectivement environ 0,69 et 0,49. Cela confirme ce qui avait déjà été observé dans l'analyse par variété :
le corse, de par sa proximité étymologique avec le français et sa structure romane, génère des paires
parallèles dont la similarité sémantique vectorielle est plus élevée que pour les autres langues. La paire
fr-oc arrive en deuxième position avec un score parallèle d'environ 0,63 et un score semi-parallèle
d'environ 0,45, ce qui est cohérent avec la nature romane de l'occitan et sa parenté avec le français.
La paire fr-gsw présente les scores les plus bas pour la classe parallèle, aux alentours de 0,49, ce qui
reflète la distance structurelle plus importante entre le français et l'alsacien, cette dernière étant une
langue germanique dont les représentations vectorielles dans l'espace sémantique multilingue de SBERT
sont naturellement moins proches de celles du français que ne le sont les langues romanes.

Les scores non parallèles sont en revanche remarquablement stables entre les trois paires de langues,
oscillant autour de 0,18 à 0,19. Cette stabilité est attendue : les paires non parallèles, construites
par appariement aléatoire de phrases, ne présentent aucun lien sémantique systématique quelle que soit
la langue régionale impliquée. L'écart entre le score parallèle et le score non parallèle, qui constitue
la marge de discrimination effective du modèle, est donc plus favorable pour le corse (0,51) que pour
l'occitan (0,45) ou l'alsacien (0,30), ce dernier présentant la marge la plus réduite et donc la
classification la plus difficile.

\section{Entraînement du modèle de régression logistique}

Une fois le jeu de données constitué et les variables d'apprentissage calculées, il devient possible
d'entraîner un modèle supervisé capable d'identifier automatiquement les paires de phrases parallèles ou
semi-parallèles au sein d'un corpus comparable non aligné. Cette étape repose sur un modèle de régression
logistique entraîné à partir du fichier produit par le notebook précédent, qui associe à chaque paire de
phrases l'ensemble des scores calculés : similarité sémantique SBERT, score de longueur, indices de Jaccard
sur les n-grammes de caractères et indice de Jaccard phonétique.

Avant l'entraînement proprement dit, une phase exploratoire permet de vérifier la cohérence du jeu de données.
Les statistiques descriptives calculées par classe confirment que les variables discriminent bien les trois
catégories : les paires parallèles obtiennent systématiquement les scores SBERT les plus élevés, suivies
des paires semi-parallèles, puis des paires non parallèles. Cette hiérarchie, observable également dans
les graphiques produits lors de la constitution du corpus, valide le bien-fondé des variables retenues pour
l'apprentissage. La matrice de corrélation entre variables permet quant à elle de détecter d'éventuelles
redondances entre les indices de Jaccard, qui mesurent des chevauchements à des granularités différentes
et peuvent donc présenter une corrélation non négligeable entre eux.

Les données sont ensuite divisées en un ensemble d'entraînement et un ensemble de test selon un découpage
stratifié de 80/20, garantissant que les trois classes sont représentées de manière proportionnelle dans
chacun des deux ensembles. Les variables sont normalisées par standardisation, une étape indispensable pour
la régression logistique dont les coefficients sont sensibles à l'échelle des variables d'entrée. Le scaler
ajusté sur les données d'entraînement est sauvegardé afin de pouvoir être appliqué de manière cohérente lors
de l'inférence sur de nouvelles données.

Deux modèles distincts sont entraînés. Le premier est un modèle binaire qui oppose les paires parallèles
et semi-parallèles aux paires non parallèles. Ce modèle constitue le cœur opérationnel du système : appliqué
à un corpus comparable, il permet de filtrer les paires présentant un degré suffisant de parallélisme pour
être exploitables. Le second est un modèle multiclasse à trois étiquettes, qui distingue les paires
parallèles, semi-parallèles et non parallèles. Ce modèle plus fin permet, pour les paires retenues par
le modèle binaire, de distinguer les correspondances exactes des correspondances partielles, une information
utile pour orienter leur usage ultérieur dans l'entraînement de systèmes de traduction automatique. Les
deux modèles utilisent une pondération équilibrée des classes afin de pallier d'éventuels déséquilibres
résiduels dans le jeu de données, et sont évalués par validation croisée à cinq plis sur les données
d'entraînement avant leur application sur le jeu de test.

L'évaluation du modèle binaire repose sur un rapport de classification complet, incluant précision, rappel
et score F1 pour chaque classe, ainsi que sur une matrice de confusion permettant d'identifier les types
d'erreurs les plus fréquents. La courbe ROC et la courbe précision-rappel complètent cette évaluation en
quantifiant la capacité du modèle à séparer les deux classes à différents seuils de décision. Ces deux
courbes sont particulièrement pertinentes dans notre contexte, où le coût d'un faux positif (une paire
non parallèle incorrectement retenue) et celui d'un faux négatif (une paire parallèle incorrectement
rejetée) ne sont pas nécessairement symétriques selon l'usage visé. L'analyse des coefficients du modèle
permet enfin d'identifier quelles variables contribuent le plus à la décision de classification. On
s'attend à ce que le score SBERT soit le contributeur principal, les variables de surface venant en
complément pour les cas où la similarité sémantique seule ne suffit pas à trancher.

Une analyse complémentaire ventile les performances du modèle binaire par langue régionale, en distinguant
le corse, l'alsacien et l'occitan. Cette désagrégation est importante dans notre contexte, car les trois
langues présentent des profils très différents en termes de volume de données d'entraînement et de proximité
avec le français. On s'attend notamment à ce que les performances sur le corse soient plus fragiles, cette
langue étant sous-représentée dans le corpus d'entraînement parallèle comme l'a montré l'analyse quantitative
des données.

Enfin, le notebook propose une fonction d'inférence réutilisable permettant d'appliquer le modèle entraîné
à de nouvelles paires de phrases sans nécessiter de rechargement du code de calcul des variables. Cette
fonction calcule l'ensemble des variables pour une paire donnée, les normalise à l'aide du scaler
sauvegardé, puis retourne la prédiction du modèle accompagnée d'un score de confiance. Elle peut être
appliquée à l'ensemble des combinaisons de phrases d'un corpus comparable non aligné, produisant une
liste de paires candidates classées par degré de parallélisme. Cette étape constitue le point
d'articulation entre le travail de classification supervisée réalisé sur les données d'entraînement
et l'application opérationnelle du modèle sur le corpus Wikimédia utilisé pour la validation externe.

\newpage

\section{Résultats de l'entraînement}


\begin{figure}[h]
    \centering
    \includegraphics[width=\textwidth]{sbert_distribution.png}
    \caption{Distribution des scores SBERT par classe (parallèles, semi et non parallèles)}
    \label{sbert_distribution}
\end{figure}

/! pour le boxplot ajouté le nombre de paires verifiées

Ce double graphique présente la distribution des scores SBERT selon les trois catégories de parallélisme.
L'histogramme de gauche révèle une structure trimodale cohérente avec les attentes théoriques. Les paires
non parallèles se concentrent massivement dans les valeurs basses, entre 0 et 0,3, avec un pic prononcé
autour de 0,05 à 0,10, reflétant l'absence de lien sémantique entre des phrases appariées aléatoirement.
Les paires semi-parallèles présentent une distribution plus étalée, centrée autour de 0,3 à 0,5,
correspondant à la troncature partielle des phrases régionales qui conserve une proximité sémantique avec
le segment français mais réduit le contenu informationnel. Les paires parallèles se distinguent nettement
avec une concentration très marquée dans les valeurs hautes, proche de 1,0, confirmant que le modèle SBERT
capture efficacement la correspondance sémantique entre traductions exactes même en contexte multilingue
régional. Le graphe de boîtes à moustaches représente les données pour 3435 paires de phrases
(soit 1145 par classe). Le boxplot de droite synthétise ces observations : la médiane des paires
non parallèles se situe autour de 0,17, celle des paires semi-parallèles autour de 0,40, et celle des
paires parallèles autour de 0,60. On notera la présence de nombreux points aberrants au-dessus de 0,55 pour
la classe non parallèle, correspondant vraisemblablement à des phrases thématiquement proches appariées
par hasard lors de la construction du sous-ensemble négatif, notamment au sein d'un même texte narratif
comme la fable d'Ésope. La variance importante de la classe parallèle, visible dans l'amplitude du boxplot,
reflète quant à elle l'hétérogénéité dialectale du corpus, les variétés LIMSI présentant des scores nettement
inférieurs aux textes ParCoLab.

\newpage

\begin{figure}[h]
    \centering
    \includegraphics[width=\textwidth]{feature_correlation.png}
    \caption{feature correlation}
    \label{feature_correlation}
\end{figure}

La matrice de corrélation entre variables d'apprentissage révèle plusieurs relations structurelles
importantes. La corrélation négative entre le score SBERT et le score de longueur (-0,58) est attendue et
cohérente : une asymétrie de longueur importante entre les deux phrases tend à signaler une correspondance
imparfaite, ce qui se traduit simultanément par un score sémantique plus faible. Entre les indices de Jaccard,
on observe une corrélation très élevée entre les trigrammes et les quadrigrammes (0,97), ce qui indique une
forte redondance entre ces deux variables. Ces deux mesures capturent des chevauchements de surface à des
granularités très proches et apportent donc une information quasi identique au modèle. La corrélation entre
l'indice de Jaccard phonétique et les trigrammes (0,76) et quadrigrammes (0,71) est également notable,
suggérant que les similarités phonétiques et les similarités de surface en caractères covarient fortement
dans notre corpus, probablement parce que les langues régionales romanes partagent avec le français des
séquences de caractères proches de leur représentation phonétique. Le score SBERT présente des corrélations
modérées et positives avec les indices de Jaccard (entre 0,35 et 0,47), indiquant que les deux familles de
variables mesurent des aspects complémentaires mais partiellement liés de la similarité entre phrases. Cette
structure de corrélation justifie a posteriori le choix de conserver l'ensemble des variables malgré les
redondances partielles, la régression logistique étant capable de pondérer leur contribution respective.

\newpage

\begin{figure}[h]
    \centering
    \includegraphics[width=\textwidth]{feature_importance.png}
    \caption{feature importance}
    \label{feature_importance}
\end{figure}

Le graphique des coefficients du modèle binaire apporte un résultat particulièrement intéressant et en partie
contre-intuitif. L'indice de Jaccard sur les quadrigrammes de caractères présente de loin le coefficient le
plus élevé (environ 3,0), dépassant significativement le score SBERT (environ 1,4) et l'indice phonétique
(environ 1,0). Ce résultat suggère que le chevauchement de séquences de quatre caractères consécutifs
constitue le signal le plus discriminant pour identifier les paires parallèles dans notre corpus, davantage
encore que la similarité sémantique vectorielle. Cela s'explique par la nature des langues étudiées{ :}
l'occitan, le corse et l'alsacien partagent avec le français de nombreuses racines lexicales dont la graphie
est suffisamment proche pour que les quadrigrammes de caractères constituent un indicateur robuste de
correspondance. Le score SBERT arrive en deuxième position, confirmant son rôle central mais montrant qu'il
n'est pas le seul déterminant fiable dans ce contexte. L'indice phonétique contribue positivement mais dans
une moindre mesure, tandis que les bigrammes de caractères ont un effet positif faible. Le coefficient négatif
des trigrammes (environ -0,4) est en apparence paradoxal mais s'explique par la forte colinéarité avec les
quadrigrammes observée dans la matrice de corrélation : le modèle compense partiellement la redondance entre
ces deux variables en attribuant un poids négatif à l'une d'elles. Enfin, le score de longueur présente un
coefficient très proche de zéro, indiquant qu'une fois les autres variables contrôlées, l'asymétrie de
longueur n'apporte qu'une contribution marginale à la classification.

\begin{figure}[h]
    \centering
    \includegraphics[width=\textwidth]{confusion_binary.png}
    \caption{confusion binary}
    \label{confusion_binary}
\end{figure}

La matrice de confusion du modèle binaire montre des résultats globalement satisfaisants. Sur les 229 paires
réellement non parallèles du jeu de test, 205 sont correctement classées, soit un taux de vrais négatifs de
89,5\%. Les 24 faux positifs restants correspondent à des paires non parallèles incorrectement identifiées
comme parallèles, probablement des phrases thématiquement proches dont les variables de surface ou sémantiques
dépassent le seuil de décision. Sur les 458 paires réellement parallèles ou semi-parallèles, 334 sont
correctement identifiées, soit un taux de vrais positifs de 73,0\%, tandis que 124 sont incorrectement
rejetées comme non parallèles. Ce taux de faux négatifs plus élevé s'explique en partie par la présence de
paires semi-parallèles dans la classe positive : leurs scores sémantiques et de surface étant intermédiaires
par construction, elles sont naturellement plus difficiles à distinguer des paires non parallèles. Ce
déséquilibre entre le taux de vrais négatifs et le taux de vrais positifs indique que le modèle est légèrement
conservateur, préférant rejeter des paires douteuses plutôt que d'introduire du bruit dans les données
extraites. Ce comportement est acceptable dans notre contexte, où la qualité des paires retenues prime sur
leur exhaustivité.

\begin{figure}[h]
    \centering
    \includegraphics[width=\textwidth]{confusion_multiclass.png}
    \caption{confusion multiclass}
    \label{confusion_multiclass}
\end{figure}

La matrice de confusion du modèle multiclasse révèle une asymétrie prononcée dans les performances selon les
classes. Les paires parallèles sont correctement classées dans 208 cas sur 229, soit un taux de 90,8\%,
ce qui confirme que le modèle identifie bien les correspondances exactes. Les paires non parallèles
obtiennent un taux de classification correcte de 78,2\% (179 sur 229), les erreurs se répartissant de
manière quasi égale entre des confusions avec la classe parallèle (24 cas) et la classe semi-parallèle
(26 cas). Les paires semi-parallèles constituent sans surprise la classe la plus difficile à identifier
correctement : seulement 114 sur 229 sont bien classées, soit 49,8\%. Les erreurs se concentrent massivement
vers la classe non parallèle (77 cas), ce qui s'explique par le fait que les phrases tronquées
artificiellement présentent des scores de surface et sémantiques qui peuvent se rapprocher de ceux des
paires non parallèles, en particulier lorsque la troncature est importante. Les 38 confusions entre
semi-parallèle et parallèle correspondent probablement à des cas où la troncature était mineure et n'a
pas suffisamment dégradé les variables pour distinguer les deux catégories. Ces résultats montrent que la
frontière entre paires semi-parallèles et non parallèles est la plus difficile à établir, ce qui est cohérent
avec la nature artificielle des semi-parallèles construites par troncature.

\newpage

\begin{figure}[h]
    \centering
    \includegraphics[width=\textwidth]{roc_pr_binary.png}
    \caption{roc pr binary}
    \label{roc_pr_binary}
\end{figure}

Les courbes ROC et précision-rappel du modèle binaire confirment de bonnes performances globales.
L'AUC-ROC de 0,867 indique que le modèle est capable de séparer correctement les paires parallèles des
paires non parallèles dans 86,7\% des cas, quelle que soit la valeur du seuil de décision. La courbe ROC
présente une montée initiale rapide, signifiant que le modèle atteint un taux de vrais positifs élevé à
faible coût en faux positifs, ce qui est particulièrement souhaitable dans notre contexte où l'on cherche
à maximiser la précision des paires extraites. La courbe précision-rappel, avec un score de précision
moyenne de 0,934, est encore plus favorable. Elle montre que pour des valeurs de rappel allant jusqu'à
environ 0,5, la précision se maintient proche de 1,0, indiquant que le modèle peut être utilisé en mode
conservateur pour extraire un sous-ensemble de paires de très haute qualité avec très peu de bruit.
La dégradation de la précision n'intervient qu'au-delà d'un rappel d'environ 0,5, lorsque le modèle
commence à récupérer les paires les plus ambiguës. Ces résultats suggèrent que l'ajustement du seuil de
décision constitue un levier important pour adapter le comportement du classifieur selon que l'on privilégie
la précision ou le rappel dans l'application finale.

\section{Validation du modèle}

Afin d’évaluer la qualité du modèle obtenu, nous procédons à une validation externe en utilisant le corpus
de \citet{tsolakis_2024} issu de données Wikimedia. Il contient des paires de phrases alignées en corse, occitan,
alsacien et français. Il présente deux caractéristiques importantes pour notre évaluation. D'une part, il est
structurellement distinct du corpus d'entraînement, reposant sur des textes encyclopédiques hétérogènes
plutôt que sur des textes narratifs ou littéraires. D'autre part, il couvre les mêmes paires de langues que
notre corpus d'entraînement, ce qui permet une évaluation directe de la généralisation du modèle.

L’application du modèle sur ce corpus permet de vérifier{ :}
\begin{itemize}
    \item la capacité du classifieur à identifier correctement des phrases parallèles dans un contexte
    différent
    \item l’absence de surapprentissage spécifique au corpus LIMSI et à ses textes majoritairement occitan
    \item la robustesse du modèle face au différence structurelle et thématique.
\end{itemize}

La phase de validation constitue l'étape finale du pipeline, permettant d'évaluer la capacité du modèle
à identifier des fragments parallèles dans un corpus comparable réel, structurellement et thématiquement
distinct du corpus d'entraînement. Elle repose sur deux programmes distincts dont la séparation est
intentionnelle : un script de préparation des données d'une part, et un notebook d'inférence et d'analyse
d'autre part.

\subsection{Préparation du corpus de validation}

Le script de préparation transforme les fichiers textes bruts issus des dossiers de validation en tableaux
structurés de paires de phrases exploitables par le classifieur. Ce découpage en deux étapes présente un
avantage pratique important : le nettoyage et la segmentation, potentiellement coûteux en temps de calcul
sur de grands corpus, ne sont effectués qu'une seule fois. Les relances ultérieures du notebook de validation
repartent directement des fichiers CSV ainsi produits, sans nécessiter de retraitement des textes sources.

Le script parcourt cinq dossiers distincts correspondant aux variétés linguistiques étudiées, à savoir
l'alsacien standard, le mulhousien et le strasbourgeois pour les variétés alsaciennes, ainsi que le corse
et l'occitan. Pour chaque dossier, les fichiers français et régionaux sont appariés par préfixe de nom
de fichier commun, garantissant que les combinaisons générées restent intra-documentaires, c'est-à-dire
que seules des phrases issues d'articles traitant du même sujet sont mises en regard. Cette contrainte
est essentielle dans le cadre d'un corpus comparable : comparer des phrases issues de documents sans
lien thématique produirait une proportion écrasante de paires non parallèles et rendrait l'extraction de
fragments pertinents impossible.

Les textes bruts présentent une structure hétérogène caractéristique des extraits encyclopédiques :
paragraphes narratifs, titres de sections, listes d'éléments, noms propres isolés et données tabulaires.
Un pipeline de nettoyage en deux étapes est donc appliqué avant toute segmentation. La première étape
supprime les lignes non informatives : les titres courts de quatre mots ou moins sans ponctuation finale,
les éléments de liste réduits à un nom propre isolé, et les lignes constituées d'une simple puce ou d'un
tiret sans contenu réel. Ces éléments, s'ils étaient conservés, produiraient des segments trop courts ou
trop peu informatifs pour être comparés de manière fiable. La seconde étape normalise les espaces multiples
et les tirets longs afin d'homogénéiser la ponctuation avant la segmentation.

La segmentation en phrases repose sur un découpage aux ponctuations fortes, c'est-à-dire le point, le
point d'exclamation et le point d'interrogation, suivies d'une majuscule. Afin d'éviter les coupures
erronées sur des abréviations courantes ou des nombres décimaux, un mécanisme de protection par
substitution temporaire est appliqué avant le découpage : les abréviations identifiées sont remplacées
par des marqueurs neutres, les nombres décimaux sont protégés de la même manière, puis ces substitutions
sont annulées après la segmentation. Les phrases de moins de cinq tokens sont ensuite écartées, car elles
ne portent pas suffisamment d'information sémantique ou lexicale pour alimenter les variables d'apprentissage
de manière fiable.

Une fois les phrases extraites, le script génère l'ensemble des combinaisons cartésiennes entre phrases
françaises et phrases régionales au sein de chaque paire de documents. Le produit cartésien, bien que
potentiellement volumineux pour des documents longs, est la stratégie naturelle pour un corpus comparable
non aligné : en l'absence d'alignement préalable, toutes les combinaisons doivent être évaluées par le
classifieur afin d'identifier celles qui présentent un degré suffisant de parallélisme. Le résultat est
exporté sous forme de CSV structuré, avec une ligne par paire de phrases et des colonnes indiquant la
variété linguistique, le code langue, la paire de documents d'origine, la phrase française et la phrase
régionale.

\subsection{Inférence et analyse}

Le notebook de validation charge les CSV produits par le script de préparation, calcule les variables
d'apprentissage sur l'ensemble des paires, puis applique les deux modèles entraînés. Le calcul des variables
suit exactement le même pipeline que celui utilisé lors de la constitution du jeu de données d'entraînement :
score de similarité sémantique SBERT, score de longueur, indices de Jaccard sur les bigrammes, trigrammes et
quadrigrammes de caractères, et indice de Jaccard phonétique. Cette cohérence est indispensable pour garantir
que les données présentées au classifieur lors de la validation sont comparables à celles sur lesquelles il
a été entraîné. Les variables sont ensuite normalisées à l'aide du scaler ajusté lors de l'entraînement, qui
est rechargé depuis le fichier sauvegardé.

L'inférence produit pour chaque paire une prédiction binaire opposant les paires parallèles ou semi-parallèles
aux paires non parallèles, accompagnée d'un score de confiance, ainsi qu'une prédiction multiclasse
distinguant les trois catégories. Le score de confiance du modèle binaire est utilisé comme critère de
filtrage : seules les paires dont la confiance dépasse un seuil paramétrable, fixé par défaut à 0,65, sont
retenues comme candidates. Ce seuil constitue le principal levier d'ajustement du comportement du système,
permettant de faire varier le compromis entre précision et rappel selon les besoins de l'application.

L'analyse des résultats est conduite à deux niveaux de granularité. Au niveau des variétés, les statistiques
permettent de comparer les performances du modèle entre les trois variétés alsaciennes d'une part, et entre
l'alsacien, le corse et l'occitan d'autre part. Cette désagrégation est particulièrement pertinente dans
notre contexte, car les variétés mulhousienne et strasbourgeoise de l'alsacien peuvent présenter des graphies
sensiblement différentes de la variété standard, ce qui est susceptible d'affecter les variables de surface
telles que les indices de Jaccard. Au niveau des langues, la comparaison permet d'évaluer dans quelle
mesure le modèle, entraîné principalement sur des données occitanes, généralise correctement à des langues
structurellement différentes comme le corse ou l'alsacien.

Les visualisations produites complètent cette analyse quantitative. La distribution des scores SBERT par
variété permet d'observer si les paires candidates présentent des scores significativement supérieurs à
l'ensemble des comparaisons effectuées, ce qui constituerait une validation indirecte de la pertinence du
filtrage. La courbe d'effet du seuil de confiance permet quant à elle d'identifier le point d'inflexion à
partir duquel une élévation du seuil entraîne une réduction rapide du nombre de candidates sans gain
substantiel en qualité moyenne, informant ainsi le choix du seuil opérationnel.

Enfin, les paires candidates retenues sont exportées sous forme de CSV, à la fois dans un fichier global
et dans des fichiers individuels par variété. Ces fichiers constituent la sortie opérationnelle du système :
ils contiennent les fragments parallèles et semi-parallèles identifiés automatiquement dans le corpus
comparable, accompagnés de l'ensemble des scores calculés, et peuvent être utilisés directement pour enrichir
les ressources parallèles disponibles pour les langues régionales étudiées ou pour alimenter l'entraînement
de systèmes de traduction automatique neuronale.

\section{Résultats de la validation}

\begin{figure}[h]
    \centering
    \includegraphics[width=\textwidth]{heatmap_sbert_validation.png}
    \caption{heatmap sbert validation}
    \label{heatmap_sbert_validation}
\end{figure}

Cette heatmap présente les scores SBERT moyens des paires candidates retenues, ventilés par variété
linguistique et par classe prédite par le modèle multiclasse. Les classes sont encodées numériquement :
0 correspond aux paires non parallèles, 1 aux paires parallèles et 2 aux paires semi-parallèles. Plusieurs
observations méritent attention.

Pour le corse et le strasbourgeois, la classe 0 affiche un score de 0,00, ce qui signifie qu'aucune paire
non parallèle n'a été retenue parmi les candidates pour ces deux variétés au seuil de confiance appliqué.
Ce résultat est cohérent avec le comportement attendu du filtre : le modèle a correctement écarté les paires
sans lien sémantique pour ces langues. Pour l'alsacien standard, le mulhousien et l'occitan en revanche, un
score non nul pour la classe 0 indique qu'un certain nombre de paires non parallèles ont tout de même
franchi le seuil de confiance, ce qui constitue des faux positifs résiduels. Ces erreurs sont
particulièrement attendues pour l'alsacien dont les trois variétés partagent avec le français des structures
lexicales de surface suffisamment proches pour induire le modèle en erreur sur des variables comme les
indices de Jaccard.

Les scores des classes parallèle et semi-parallèle sont remarquablement stables entre variétés, oscillant
entre 0,41 et 0,54. Le corse présente les scores les plus élevés pour ces deux classes (0,54), ce qui
suggère que les paires candidates retenues pour cette langue présentent une proximité sémantique plus
marquée que pour les autres variétés, probablement en raison de la plus grande proximité structurelle du
corse avec le français comparativement à l'alsacien. L'absence de différence notable entre les scores des
classes parallèle et semi-parallèle pour la plupart des variétés confirme la difficulté, déjà observée
lors de l'évaluation sur le jeu de test, à séparer ces deux catégories par le seul score sémantique.

\newpage

\begin{figure}[h]
    \centering
    \includegraphics[width=\textwidth]{Sbert_distribution_by_lp.png}
    \caption{Sbert distribution by lp}
    \label{Sbert_distribution_by_lp}
\end{figure}

Ce graphique présente en cinq histogrammes la distribution des scores SBERT pour chaque variété,
superposant en bleu l'ensemble des comparaisons effectuées et en vert les seules paires candidates retenues.
La ligne pointillée verticale indique le score SBERT moyen des candidates.

La première observation transversale à toutes les variétés est la concentration massive des comparaisons
totales dans les valeurs très basses, entre 0 et 0,2. Ce phénomène est inhérent à la stratégie de produit
cartésien : la très grande majorité des combinaisons de phrases générées au sein d'un même document
comparable sont sémantiquement éloignées, ce qui confirme que le corpus utilisé est bien un corpus
comparable et non un corpus parallèle. Le volume de comparaisons varie considérablement selon les variétés,
les histogrammes de l'alsacien standard et de l'occitan atteignant respectivement environ 140 000 et
260 000 comparaisons, reflétant la plus grande taille des corpus disponibles pour ces variétés.

Les paires candidates, représentées en vert, apparaissent systématiquement dans les valeurs intermédiaires
à hautes du score SBERT, avec des moyennes comprises entre 0,42 et 0,54 selon les variétés. Pour le corse
et le strasbourgeois, dont les corpus sont plus petits, les candidates sont peu nombreuses mais bien
concentrées au-delà de 0,4, indiquant que le filtre a correctement ciblé les paires les plus sémantiquement
proches. Pour l'occitan et les variétés alsaciennes, la distribution des candidates est plus étalée, ce qui
s'explique par le volume plus important de comparaisons disponibles et donc une plus grande diversité de
paires potentiellement ambiguës atteignant le seuil de confiance.

\begin{figure}[h]
    \centering
    \includegraphics[width=\textwidth]{trust_threshold.png}
    \caption{Seuil de confiance par sbert moyen et nombre de candidats}
    \label{trust_threshold}
\end{figure}

Ce graphique à double axe permet d'observer simultanément l'évolution du nombre de paires candidates et
du score SBERT moyen de ces candidates en fonction du seuil de confiance appliqué, le seuil actuel de 0,65
étant matérialisé par la ligne pointillée rouge.

Le nombre de paires candidates décroît de manière monotone et régulière à mesure que le seuil augmente,
passant d'environ 280 000 paires à un seuil de 0,50 à environ 25 000 paires à un seuil de 0,95. La
décroissance est approximativement linéaire entre 0,50 et 0,75, puis ralentit légèrement au-delà,
indiquant qu'un noyau de paires bénéficie d'une confiance très élevée du modèle. Au seuil actuel de 0,65,
environ 160 000 paires sont retenues, ce qui représente un volume conséquent mais cohérent avec la taille
totale du corpus comparable traité.

Le score SBERT moyen des candidates évolue en sens inverse, progressant de 0,38 à un seuil de 0,50 jusqu'à
environ 0,63 à un seuil de 0,95. Cette progression confirme que l'élévation du seuil de confiance permet
effectivement de sélectionner des paires présentant une proximité sémantique plus élevée. On observe
cependant que la progression du score SBERT moyen s'accélère nettement au-delà de 0,75, passant de 0,49
à 0,63 entre 0,75 et 0,95, alors que le gain entre 0,50 et 0,75 est plus modeste (de 0,38 à 0,49). Ce point
d'inflexion autour de 0,75 constitue une information utile pour le réglage du seuil : un seuil dans cette
zone permettrait de réduire significativement le volume de candidates tout en garantissant une qualité
sémantique nettement supérieure à celle obtenue avec le seuil actuel de 0,65.